\chapter{Conclusão}
\label{chap:conclusion}

\section{Considerações Finais}
\label{sec:final-considerations}

Este trabalho apresentou o desenvolvimento da Medicano,
uma plataforma web de agendamento de serviços de saúde
com triagem inteligente por chatbot, construída como
trabalho de conclusão de curso do Bacharelado em
Engenharia de Software da PUC Campinas.

O problema endereçado é concreto e relevante: 91\% dos
centros médicos brasileiros ainda utilizam o telefone
como principal canal de agendamento
\cite{doctoralia2022}, e uma parcela significativa dos
pacientes não sabe qual especialidade procurar diante
de seus sintomas. A Medicano propõe uma solução que
atua simultaneamente nas duas frentes — a digitalização
do agendamento e a orientação inteligente ao paciente
por meio de um chatbot baseado em modelos de linguagem
de grande escala.

Ao longo do primeiro semestre letivo de 2026, a equipe
percorreu sete sprints quinzenais, cobrindo desde a
configuração da infraestrutura e a autenticação até o
módulo de agendamentos, o painel da clínica, o módulo
de assinaturas e o deploy em produção na AWS. O
resultado é um sistema funcional, acessível pelo
domínio \texttt{medicano.app}, com frontend distribuído
via Amazon CloudFront e backend exposto por Nginx em
instância Amazon EC2.

A adoção do \textit{Spec Driven Development} como
metodologia de especificação, complementada pelo Scrum
para execução iterativa, mostrou-se eficaz para um
projeto com cinco desenvolvedores simultâneos: a
exigência de spec aprovada antes de cada sprint reduziu
o retrabalho e permitiu que diferentes membros da
equipe implementassem módulos em paralelo a partir de
contratos previamente acordados.

\section{Contribuições do Trabalho}
\label{sec:contributions}

As principais contribuições técnicas deste trabalho são:

\begin{itemize}
  \item \textbf{Chatbot de triagem com LLM:} integração
        do Vercel AI SDK com modelos de linguagem para
        orientar o paciente à especialidade mais
        adequada em linguagem natural, com respostas
        em modo \textit{streaming} para experiência
        conversacional fluida

  \item \textbf{Plataforma multilocatária:} arquitetura
        que suporta simultaneamente clínicas com acesso
        delegado a funcionários, profissionais autônomos
        com agenda individual e pacientes sem vínculo
        institucional, mantendo isolamento de dados
        entre os perfis

  \item \textbf{Gestão de sessões com revogação
        imediata:} uso do Redis para armazenar tokens
        JWT ativos com TTL automático, permitindo
        invalidação instantânea de sessão no logout
        sem dependência exclusiva da assinatura
        criptográfica do token

  \item \textbf{Infraestrutura reproduzível por
        ambiente:} estratégia de databases separados
        no mesmo cluster Atlas M0 para desenvolvimento,
        \textit{staging} e produção, com secrets
        gerenciados pelo AWS Secrets Manager e
        diferenciados por prefixo — eliminando o risco
        de credenciais versionadas e garantindo
        paridade de configuração entre ambientes

  \item \textbf{Documentação orientada por spec:}
        aplicação do SDD em um projeto acadêmico de
        múltiplos autores por meio de \textit{specs}
        versionadas no próprio repositório, servindo
        diretamente de entrada para o \textit{pipeline}
        de desenvolvimento assistido por IA e
        demonstrando a viabilidade da abordagem para
        equipes pequenas com alto grau de paralelismo
\end{itemize}

\section{Limitações}
\label{sec:limitations}

O escopo do presente trabalho impõe algumas limitações
que devem ser consideradas na avaliação do sistema:

\begin{itemize}
  \item \textbf{Capacidade do cluster Atlas M0:}
        o plano gratuito do MongoDB Atlas impõe limites
        de armazenamento e desempenho que podem se
        tornar restritivos com crescimento do volume
        de dados em produção

  \item \textbf{Ausência de aplicativo móvel:}
        a plataforma é integralmente web-based, sem
        aplicativo nativo para iOS ou Android, o que
        pode limitar a experiência em dispositivos
        móveis

  \item \textbf{Triagem restrita ao direcionamento
        de especialidade:} o chatbot não realiza
        diagnóstico clínico — direciona o paciente
        à especialidade adequada, mas não substitui
        a avaliação médica. Esse limite é
        intencional e comunicado ao usuário na
        interface

  \item \textbf{Integração de pagamento não incluída:}
        o módulo de assinaturas implementado no
        primeiro semestre cobre a estrutura de planos
        e controle de acesso, mas a integração com
        gateway de pagamento está planejada para o
        segundo semestre

  \item \textbf{Cobertura de testes concentrada nos
        fluxos críticos:} testes automatizados cobrem
        autenticação, regras de agendamento e triagem;
        módulos secundários dependem de validação
        manual ao final de cada sprint
\end{itemize}

\section{Trabalhos Futuros}
\label{sec:future-work}

O segundo semestre letivo de 2026 será dedicado à
expansão de escopo da plataforma, cujo levantamento
de requisitos ocorrerá no início do período. Além
das funcionalidades a serem definidas nessa etapa,
algumas evoluções técnicas já identificadas incluem:

\begin{itemize}
  \item \textbf{Integração com gateway de pagamento:}
        conclusão do módulo de assinaturas com
        cobrança recorrente, permitindo que clínicas
        e profissionais autônomos ativem seus planos
        diretamente pela plataforma

  \item \textbf{Aplicativo móvel:} desenvolvimento
        de cliente React Native para iOS e Android,
        reutilizando os tipos compartilhados do
        monorepo e consumindo a mesma API existente

  \item \textbf{Integração com a RNDS:} conexão com
        a Rede Nacional de Dados em Saúde do
        Ministério da Saúde, habilitando o acesso ao
        histórico clínico do paciente no contexto
        do agendamento

  \item \textbf{Telemedicina:} adição de fluxo de
        consulta por videochamada integrado ao
        agendamento, com salas virtuais gerenciadas
        pela plataforma

  \item \textbf{Escalabilidade horizontal:}
        migração do Redis para Amazon ElastiCache e
        do backend para Amazon ECS, habilitando
        múltiplas instâncias do servidor sem
        dependência de estado local

  \item \textbf{Extração do módulo de triagem:}
        separação do chatbot de triagem em um
        microsserviço independente, permitindo
        evoluir a lógica de LLM sem impacto no
        backend principal e habilitando o uso de
        modelos especializados em saúde
\end{itemize}